\documentclass[12pt]{article}
\usepackage[T1]{fontenc}
\usepackage[utf8]{inputenc}
\usepackage{lmodern}
\usepackage{textcomp}
\usepackage{enumitem}
\usepackage{geometry}
\geometry{margin=1in}
\begin{document}
\begin{center}
Answer Key - Science - K-12 Level Assessment 11 \
\small (Answers are ordered exactly as in the question paper.)
\end{center}

\section*{Multiple Choice}
\begin{enumerate}[label=\arabic*.]
\item \textbf{Answer:} A
\\ \textit{Options:} A) heavy metals threaten some forms of life.; B) heavy metals will increase available resources.; C) exposure to heavy metals causes healthy mutations.; D) mining for heavy metals causes environmental stability.
\\ \textit{Note:} Heavy metals like lead and mercury are toxic and can accumulate in living organisms, leading to harmful effects on health and biodiversity, which is why scientists are concerned about their presence in the environment.
\item \textbf{Answer:} C
\\ \textit{Options:} A) aluminum; B) copper; C) plastic; D) water
\\ \textit{Note:} Plastic is an insulator, meaning it does not allow electricity to flow through it easily, in contrast to conductive materials like metals.
\item \textbf{Answer:} B
\\ \textit{Options:} A) conservation of genetic information; B) conservation of cell resources; C) trait adaptation to environmental change; D) trait inheritance in offspring
\\ \textit{Note:} The regulation of gene expression allows cells to produce only the proteins necessary for their current functions, thereby conserving energy and resources by preventing unnecessary synthesis of proteins.
\item \textbf{Answer:} B
\\ \textit{Options:} A) making a den in the hollow of a tree; B) opening a trashcan looking for food; C) moving four legs while running; D) having a litter of four babies
\\ \textit{Note:} Opening a trashcan looking for food is a learned behavior for raccoons because it involves problem-solving and adaptability, demonstrating their ability to modify their actions based on experience in seeking food resources in human environments.
\item \textbf{Answer:} D
\\ \textit{Options:} A) seas that have evaporated.; B) natural gas subjected to heat.; C) melted and cooled metamorphic rock.; D) plant remains decomposed under pressure.
\\ \textit{Note:} Coal is formed from the remains of ancient plants that have undergone a process of decomposition and transformation under high pressure and temperature over millions of years, making this answer correct.
\end{enumerate}

\section*{True / False}
\begin{enumerate}[label=\arabic*.]
\setcounter{enumi}{5}
\item \textbf{Answer:} True
\\ \textit{Options:} A) True; B) False
\\ \textit{Note:} The statement is true because the baking process involves a chemical change where the ingredients in the dough undergo reactions that alter their molecular structure, resulting in the formation of bread.
\item \textbf{Answer:} True
\\ \textit{Options:} A) True; B) False
\\ \textit{Note:} The statement is true because the equation shows that the number of atoms of each element is the same on both sides, indicating that mass is conserved during the combustion of ethane.
\item \textbf{Answer:} False
\\ \textit{Options:} A) True; B) False
\\ \textit{Note:} Even if a compound is nontoxic, it may still disrupt water treatment processes or harm aquatic ecosystems, so it should not be poured down the drain without proper disposal methods.
\item \textbf{Answer:} True
\\ \textit{Options:} A) True; B) False
\\ \textit{Note:} Accurate records of consistent data are essential for scientists to validate their discoveries, as they provide the necessary evidence to support or refute hypotheses and ensure the reliability of experimental results.
\item \textbf{Answer:} False
\\ \textit{Options:} A) True; B) False
\\ \textit{Note:} Plants wilt on hot days primarily due to water loss through transpiration exceeding water uptake, rather than their photosynthesis rate.
\end{enumerate}

\section*{Long Form Response}
\begin{enumerate}[label=\arabic*.]
\setcounter{enumi}{10}
\item \textbf{Answer:} Chloroplasts and vacuoles play essential roles in plant cells, particularly in energy storage and utilization. Chloroplasts are the sites of photosynthesis, where they capture sunlight and convert it into chemical energy in the form of glucose. This process not only provides energy for the plant but also produces oxygen as a byproduct. Vacuoles, on the other hand, are primarily responsible for storing nutrients, waste products, and water, helping to maintain the cell's structure and pressure. By storing glucose and other compounds, vacuoles assist in energy management, allowing the plant to use this energy when needed for growth and metabolism. Together, chloroplasts and vacuoles ensure that plants can efficiently harness and utilize energy from sunlight.
\\ \textit{Note:} The answer correctly identifies chloroplasts as the organelles responsible for photosynthesis and energy conversion, while also highlighting the role of vacuoles in storing essential substances, thus illustrating the complementary functions of these structures in energy storage and cellular maintenance in plant cells.
\item \textbf{Answer:} To differentiate between table salt and table sugar, students can use several methods, such as observing their physical properties, conducting solubility tests, and performing taste tests (with caution). For example, students can note that table salt has a crystalline structure and a salty taste, while table sugar is sweet and has a different crystalline form. Accurate labeling is crucial in scientific investigations, as it ensures that substances are correctly identified, which helps avoid confusion and potential safety hazards. By using these methods, students can effectively distinguish the two substances while emphasizing the importance of proper labeling in experiments.
\\ \textit{Note:} correct because it outlines specific methods for distinguishing between table salt and table sugar, such as observing physical properties and conducting tests, while also emphasizing the significance of accurate labeling in preventing misidentification during scientific investigations.
\end{enumerate}

\vfill
\noindent\textit{End of Answer Key}
\end{document}